\documentclass{standalone}
\usepackage{tikz}

\begin{document}
\begin{tikzpicture}
    % Define coordinates to reconstruct the geometry
    % C is at the origin (0,0)
    \coordinate (C) at (0,0);
    % B is on the x-axis
    \coordinate (B) at (5,0);
    % D (now labeled P) is on the y-axis
    \coordinate (D) at (0,1.5);
    % A is positioned to create the specific shape and right angle at A
    % Slopes of DA and AB are roughly 1 and -1 to satisfy perpendicularity visually
    \coordinate (A) at (1.75,3.25);

    % Draw the quadrilateral
    \draw (C) -- (B) -- (A) -- (D) -- cycle;

    % Draw right angle mark at C
    \draw (0,0.3) -- (0.3,0.3) -- (0.3,0);

    % Draw right angle mark at A
    % Based on the coordinates, the legs have slopes 1 and -1.
    % The mark is drawn perpendicular to these legs.
    \draw (1.55,3.05) -- (1.75,2.85) -- (1.95,3.05);

    % Labels
    % Change label D to P as requested
    \node[left] at (D) {$P$};
    \node[below left] at (C) {$C$};
    \node[below right] at (B) {$B$};
    \node[above right] at (A) {$A$};
\end{tikzpicture}
\end{document}
